% Quiz 9: shortlist A, E, C, and G, in that order.
% Solutions: q9-sol-261-II2026.tex. Sources: q9-shortlist-261-II2026.md.
\documentclass[12pt]{exam}
\usepackage{amsthm}
\usepackage{bm}
\usepackage{libertine}
\usepackage[utf8]{inputenc}
\usepackage[margin=0.5in]{geometry}
\usepackage{amsmath,amssymb}
\usepackage{multicol}
\usepackage[shortlabels]{enumitem}
\usepackage{siunitx}
\usepackage{physics}
\usepackage{booktabs}
\usepackage{graphicx}
\usepackage{pgfplots}
\usepackage{listings}
\usepackage{tikz}
\usepackage{tikz-3dplot}
\usetikzlibrary{arrows.meta,decorations.markings}

\bmdefine{\ii}{i}
\bmdefine{\jj}{j}
\bmdefine{\kk}{k}
\newcommand{\te}{\theta}
\newcommand{\cC}{\mathcal{C}}
\newcommand{\word}[1]{\quad\text{#1}\quad}

\pgfplotsset{width=10cm,compat=1.9}
\usepgfplotslibrary{external}
\tikzexternalize
\tikzexternaldisable

\newcommand{\class}{Math 261-001}
\newcommand{\examnum}{Quiz 9}
\newcommand{\examdate}{November 6}

\begin{document}
\pagestyle{plain}
\thispagestyle{empty}

\noindent
\textbf{\class}\\
\textbf{\examnum}, \textbf{\examdate}

\setlength{\tabcolsep}{3.5cm}
\renewcommand{\arraystretch}{1.5}
\setlength\extrarowheight{1cm}
\begin{tabular}{ |c|c| }
 \hline
 Name   & CSU ID \#  \\
 \hline
\end{tabular}

\vspace{10pt}

Choose and answer \textbf{exactly two} of the four problems. If you attempt
more than two, only the first two will be graded. Write your final answers
in the boxes provided. Only answers in the boxes will be graded.
You may use the surrounding space for scratch work. Each problem is worth
the same number of points.

\begin{enumerate}

\item For each item state the \textbf{type of
integral} (line or surface) and the \textbf{vector field being integrated} (curl or original) after applying Stokes' theorem.
Do not compute derivatives or evaluate either integral.

\begin{enumerate}[label=(\alph*),itemsep=6pt]
\item Let $\vec F=(y,z,x)$, and let $H$ be the outward-oriented upper
hemisphere $x^2+y^2+z^2=4$, $z\geq0$. The requested integral is
\[
\iint_H(\operatorname{curl}\vec F)\cdot\vec N\,d\sigma.
\]
\item Let $\vec G=(z,x,y)$, and let $T$ be the triangle traversed in the order
\[
(2,0,0)\longrightarrow(0,2,0)\longrightarrow(0,0,2)
\longrightarrow(2,0,0).
\]
The requested integral is $\displaystyle\oint_T\vec G\cdot\vec T\,ds$.
\end{enumerate}

\begin{flushright}
\begin{minipage}{0.47\linewidth}
\centering
(a)\par\smallskip
\framebox(230,55){}
\end{minipage}\hspace{0.3cm}
\begin{minipage}{0.47\linewidth}
\centering
(b)\par\smallskip
\framebox(230,55){}
\end{minipage}
\end{flushright}

\item Let $S$ be the portion of the paraboloid \(z=x^2+y^2\) for \(1\leq z\leq 4\)
with normal pointing outward, away from the $z$-axis. Give parametrizations of its boundary curves when viewed from above. [Hint: Watch out for the induced orientation].

\vfill
\begin{flushright}
\begin{minipage}{0.47\linewidth}
\centering
Outer boundary\par\smallskip
\framebox(230,55){}
\end{minipage}\hspace{0.3cm}
\begin{minipage}{0.47\linewidth}
\centering
Inner boundary\par\smallskip
\framebox(230,55){}
\end{minipage}
\end{flushright}

\newpage

\item Find the flux of $\vec F=(y,x,y+z)$ through the portion of
$2x+y+z=2$ in the first octant, oriented \textbf{upward}: \(\iint_S\vec F\cdot\vec N\,d\sigma\).

\vfill
\begin{flushright}
\framebox(250,50){}
\end{flushright}

\item Let
\[
\vec F=(z-x^{2026},\ x-y^{2026},\ y-z^{2026}),
\]
and let $T$ be the triangle traversed in the order
\[
(2,0,0)\longrightarrow(0,2,0)\longrightarrow(0,0,2)
\longrightarrow(2,0,0).
\]
Use Stokes' theorem to evaluate
$\displaystyle\oint_T\vec F\cdot\vec T\,ds$.

\vfill
\begin{flushright}
\framebox(250,50){}
\end{flushright}

\end{enumerate}
\end{document}
